\chapter{An explicit volume theorem for ordered trees of crossing
circles}\label{ch:quantitative-circle-trees}

The preceding filling theorem makes every twist long. For the ordered
circle trees considered here, high degree can play the same role.
A degree-$d_v$ circle meets $d_v$ edge pieces. The negative second
derivatives of their local phases add at that vertex in the quantum
sum, and their disk sectors concatenate on its cusp torus.
Both calculations lead to the quantity $d_v/(d_v^2+n_v^2)$.
We obtain an explicit volume theorem allowing short or zero powers
at vertices whose degrees are sufficiently large.

\section{The knot construction}\label{sec:circle-the-knot-construction}

A two-string tangle is drawn in a ball with four ordered boundary
points: lower left, lower right, upper right, upper left. A product
places one ball above another and joins their two adjacent endpoints by
uncrossed arcs. The order of the factors is retained. The operation
\(T\) rotates an entire ball through the specified quarter-turn in the
diagram plane. Cylindrical closure joins lower left to upper left and
lower right to upper right outside the final ball.

Fix a finite tree \(\mathcal T\) with at least one edge, together with
the cyclic order of the edges and the tangle frame at each vertex. There
is one original crossing circle at every vertex, including the terminal
vertices. There is no unmarked product junction. Each edge specifies a quarter-turn attachment of a child tangle.
Crossings of edges in a drawing of the abstract tree have no meaning
for the knot diagram.

To specify the order and the frames, choose a terminal vertex as root.
At a nonroot vertex \(v\), put the integer
braid \(R^{n_v}\) between the two ordered lists of its child tangles: \begin{equation}
 W_v=\left(\prod_{w\in I_v^-}T(W_w)\right)R^{n_v}
       \left(\prod_{w\in I_v^+}T(W_w)\right).
 \label{eq:circle-1}
\end{equation} The lists \(I_v^-\) and \(I_v^+\) specify the position of the
crossing box in the original cyclic order. Empty lists are allowed. If
\(w\) is the root's unique neighbor, close \(R^{n_0}T(W_w)\)
cylindrically. The expression specifies the link diagram and its crossing disks,
not just a scalar evaluation. Other choices of cutting
section are related by the physical closed-collar moves, with their
induced rotations and reversals of factor order retained.

The augmented diagram replaces each \(R^{n_v}\) by its crossing circle
and a retained half twist \(\epsilon_v\in\{0,\pm1\}\), where
\(n_v\equiv\epsilon_v\pmod2\). Filling that circle on its original
disk-framed slope \(\mu_v+p_v\lambda_v\), with \begin{equation}
                         n_v=\epsilon_v-2p_v,
 \label{eq:circle-2}
\end{equation} restores exactly the indicated integer braid. Here \(\mu_v\) goes
once around the auxiliary circle and \(\lambda_v\) follows its punctured
disk. The main strands are not filled.

We require that the closed main strands form one knot \(K\). This
includes excluding a closed main component created inside an
intermediate box. The test uses only the strand pairings and the
parities of the \(n_v\); large magnitudes alone do not guarantee it.
Write \(d_v\) for the degree of \(v\) in the circle tree and
\(E=|E(\mathcal T)|\). This graph degree is unrelated to the quantum
level \(N\).

\section{The quantitative theorem}\label{sec:circle-the-quantitative-theorem}

Let \(F_N(M)\) be the zero-framed reduced Akutsu--Deguchi--Ohtsuki (ADO)
invariant of the knot, normalized to one on the unknot. Its
marked meridian eigenvalue is \(M\), and at \(q_N=e^{\pi i/N}\) its
complete value is the reduced Kashaev invariant \(K_N(K)=F_N(1)\).
Moving the meridian on the scale considered here means setting
\(M=q_N^s\), with \(s\) in a fixed complex disk. Define \begin{equation}
 \mathfrak S(\mathcal T,\mathbf n)
       =\sum_v\frac{d_v}{d_v^2+n_v^2},\qquad
 H_N(s)=\frac{\sin(\pi s)}{N\sin(\pi s/N)},\quad H_N(0)=1.
 \label{eq:circle-3}
\end{equation} Let \(w(D)\) be the oriented writhe of the filled diagram,
using the positive-crossing convention that defines \(R\). The exact
framing and crossing factor is \begin{equation}
 A_N(s)=q_N^{a(D)s^2/2},\qquad
                         a(D)=-w(D)-\sum_v n_v.
 \label{eq:circle-4}
\end{equation} The power uses the fixed logarithm \(\log q_N=\pi i/N\). All crossing
and strand-turn factors are included in \eqref{eq:circle-4}; it is determined by the diagram, rather than chosen after the asymptotic calculation.

\begin{theorem}[A degree and twist criterion]\label{thm:circle-quantitative} Fix the ordered tree, its diagram frames, and its
signed integer powers before letting \(N\) tend to infinity. Suppose the
main strands form a knot and \begin{equation}
                         \mathfrak S(\mathcal T,\mathbf n)\le2^{-25}.
 \label{eq:circle-5}
\end{equation} Then the knot complement is hyperbolic, and \begin{equation}
 \lim_{N\to\infty}\frac{2\pi}{N}\log|K_N(K)|
                      =\operatorname{Vol}(S^3\setminus K).
 \label{eq:circle-6}
\end{equation} For every fixed \(R<1\), the stronger statements \begin{equation}
 \frac{F_N(q_N^s)}{F_N(1)A_N(s)H_N(s)}
       =1+O_{\mathcal T,\mathbf n,R}(N^{-1}),
 \qquad |s|\le R,
 \label{eq:circle-7}
\end{equation} and \begin{equation}
 \frac{2\pi}{N}\log|F_N(q_N^s)|
           \longrightarrow\operatorname{Vol}(S^3\setminus K)
 \label{eq:circle-8}
\end{equation} hold uniformly on that disk. In particular these values are nonzero
there for all sufficiently large levels. Both level parities and
composite levels are included. The tree, its degrees, its twist numbers
and the radius \(R\) do not grow with \(N\).

\end{theorem}

The constant in \eqref{eq:circle-5} is a sufficient bound, not an
optimal threshold. The proof below uses the finite edge kernels and
their octahedral geometry directly. In the two-circle case, two odd
crossing powers give two main components rather than a knot, so that
case is excluded by the endpoint test already imposed above.

Two consequences make the range of \eqref{eq:circle-5} concrete.

\begin{itemize}

\item
  For a tree with \(E\) edges, every tuple satisfying
  \(|n_v|\ge8192\sqrt E\) satisfies \eqref{eq:circle-5}, because \(\sum_vd_v=2E\).
\item
  Start with any ordered tree on \(m\) vertices, with arbitrary fixed
  powers at those vertices. Attach \(B\) new terminal circles at each
  vertex and give each added circle a power of magnitude at least \(B\).
  The resulting sum is strictly less than \(2m/B\). Hence
  \(B\ge2^{26}m\) suffices, with the original powers still arbitrary,
  including zero. The knot/parity test remains necessary.
\end{itemize}

A product junction without a crossing circle can be included by
inserting the identity braid \(R^0\) at its prescribed position in the ordered
stack. Add a crossing circle there in the augmented diagram and fill it
with \(\epsilon=p=n=0\). The filled diagram is unchanged. The augmented
exterior is different, and the added degree-\(d\) circle contributes
\(1/d\) to \eqref{eq:circle-3}. This gives a direct application to constructions with
sufficiently large-degree unmarked junctions when the same total bound
holds. It does not treat such junctions as free or assert that a
low-degree unmarked junction satisfies \eqref{eq:circle-5}.

These constructions can be large, but their size is fixed before the
level tends to infinity. Comparisons with other volume conjectures require their invariant
normalizations and root conventions to be specified.

\section{Why the original exterior has the required geometric
pieces}\label{sec:circle-why-the-original-exterior-has-the-required-geometric-pieces}

Cut the augmented diagram along the projection sphere and the original
crossing disks. The two resulting balls carry white projection faces and
shaded disk triangles. At a leaf with half-segments
\((v_0,v_1,v_2,v_3)\), the shaded triangles are \((O_v,v_1,v_2)\) and
\((O_v,v_3,v_0)\), where \(O_v\) is the auxiliary-circle end. The
strand matching fixes every cyclic vertex order and hence every
white face boundary.

The geometric induction follows the same operation as the tangle
construction: attach one terminal child in its prescribed slot.
Let \(h_v\) denote the quarter-turned crossing-circle tangle. Write
its parent's unturned box as \(P_v=T^{-1}h_v\), so its exposed
tuple is \((v_3,v_0,v_1,v_2)\). If the old downstream matches are
\(v_1-x_1\) and \(v_2-x_2\), insertion of \(h_w\) immediately after
\(P_v\) replaces them by \begin{equation}
 v_1-w_1,\quad v_2-w_0,\quad w_2-x_1,\quad w_3-x_2.
 \label{eq:circle-9}
\end{equation} Set \(C=(v_1,x_1)\), \(B=(v_2,x_2)\), \(D=(v_1,w_1)\), and
\(A=(v_2,w_0)\). An octahedron with opposite vertices \(O_v,O_w\) and
equator \(A,B,C,D\) has shaded faces \begin{equation}
 (O_v,B,C),\quad(O_v,D,A),\quad(O_w,D,C),\quad(O_w,B,A).
 \label{eq:circle-10}
\end{equation} Glue its first face to the old face \((O_v,C,B)\), with the opposite
boundary orientation. The other three faces are precisely the new shaded
faces prescribed by \eqref{eq:circle-9}. The adjacent white sectors also concatenate in
the prescribed strand order. Thus the attachment gives the boundary of the cut diagram, with
its strand labels, before any hyperbolic shapes are assigned.

\begin{figure}[ht]
\centering
% Native TikZ fragment: a labeled ideal-cell attachment.
\begin{tikzpicture}[x=1cm,y=1cm,line cap=round,line join=round,
  font=\small,vertex/.style={circle,fill=black,inner sep=1.25pt},
  face/.style={fill=blue!15,draw=blue!65!black,line width=1pt},
  dir/.style={-{Latex[length=1.7mm]},blue!70!black,line width=.85pt}]
\coordinate (ov) at (-3.5,1.85);
\coordinate (b) at (-2.45,.65);
\coordinate (c) at (-4.65,.45);
\coordinate (l) at (-5.7,-.55);
\coordinate (ll) at (-4.85,-1.6);
\coordinate (r) at (-2.85,-1.45);
\fill[gray!7] (ov)--(b)--(r)--(ll)--(l)--(c)--cycle;
\draw[gray!65] (ov)--(b)--(r)--(ll)--(l)--(c)--cycle;
\draw[gray!55] (l)--(r) (ll)--(c) (c)--(r) (ov)--(l);
\path[face] (ov)--(c)--(b)--cycle;
\draw[dir] ($(ov)!.31!(c)$)--($(ov)!.66!(c)$);
\draw[dir] ($(c)!.34!(b)$)--($(c)!.67!(b)$);
\draw[dir] ($(b)!.30!(ov)$)--($(b)!.62!(ov)$);
\foreach \p/\lab/\where in {ov/O_v/above,b/B/right,c/C/left}{
 \node[vertex] at (\p) {};
 \node[\where=3pt] at (\p) {$\lab$};
}
\node[align=center,text width=4.1cm] at (-4.05,-2.65)
 {Existing projection ball\\shaded face $(O_v,C,B)$};
\coordinate (nv) at (2.7,1.85);
\coordinate (nw) at (2.7,-1.8);
\coordinate (na) at (3.95,-.45);
\coordinate (nb) at (3.45,.65);
\coordinate (nc) at (1.45,.45);
\coordinate (nd) at (1.95,-.65);
\fill[gray!5] (nv)--(nb)--(na)--(nw)--(nd)--(nc)--cycle;
\draw[gray!55,dashed] (nw)--(nb) (nb)--(na);
\draw[gray!70] (nv)--(na)--(nw)--(nd)--(nv)
 (nv)--(nb)--(nc)--(nv) (nc)--(nd) (nd)--(na) (nc)--(nw);
\path[face] (nv)--(nb)--(nc)--cycle;
\draw[dir] ($(nv)!.31!(nb)$)--($(nv)!.65!(nb)$);
\draw[dir] ($(nb)!.32!(nc)$)--($(nb)!.66!(nc)$);
\draw[dir] ($(nc)!.29!(nv)$)--($(nc)!.62!(nv)$);
\foreach \p/\lab/\where in {nv/O_v/above,nw/O_w/right,na/A/right,nb/B/right,nc/C/left,nd/D/left}{
 \node[vertex] at (\p) {};
 \node[\where=3pt] at (\p) {$\lab$};
}
\node[align=center,text width=4.1cm] at (2.7,-2.85)
 {\mbox{New octahedron}\\for edge $vw$\\shaded face $(O_v,B,C)$};
\node[fill=white,inner sep=3pt,align=center,text width=2.3cm]
 at (-.25,1.1) {Match equal labels;\\reverse face orientation};
\node[align=center,text width=11.2cm] at (-.85,-3.8)
 {The glued triangle becomes internal. The three other shaded faces of the
 octahedron retain the circle $v$ and introduce the circle $w$.};
\end{tikzpicture}

\caption{Attaching one octahedron in a projection half of the augmented tangle diagram. Equal labels are identified, and the blue
faces have opposite boundary orientations. The triangle becomes
internal; the other three shaded disk faces retain the parent circle
and introduce the child circle, as listed in
\eqref{eq:circle-10}. Only the attachment faces are blue. The labels
$O_v,O_w$ belong to the original auxiliary crossing circles.
The same attachment is made in the other projection half.}
\label{fig:circle-octahedral-attachment}
\end{figure}

The left-insertion version attaches along the other parent triangle. It
allows the parent box to remain at its specified position among its
children. Start from one tree edge and insert the remaining children in
the appropriate left and right order, carrying each rotated child ball
with its frame. Each projection half is a tree of \(E\) octahedra
attached along \(E-1\) triangles. Its remaining shaded faces are exactly
the original two per circle.

The internal interfaces give the additional disk pages. At a circle of
degree \(d_v\), the \(d_v\) incident blocks occur in the original cyclic
order. One distinguished seam retains the original half twist; the other
seams are the untwisted insertion maps. In particular the construction
does not put one copy of the retained crossing on every incident edge.
It also retains the half-turn produced when a turn is moved
across the closing sphere. Restoring the labeled faces recovers the
original augmented link and its primitive disk paths. Section~\ref{sec:ft} gives the full endpoint and face-map calculation.

\section{Positive cells and primitive filling periods}\label{sec:circle-a-positive-geometric-path-rather-than-a-small-total-logarithm}

One edge sum supplies the local action for both the quantum and the
geometric calculations. If \(t,u\) are its endpoint angles, that action
is \begin{equation}
 \begin{aligned}
 \psi(t,u)&=\operatorname*{stat}_{x}[b(x,t)+b(x,u)],\\
 b(x,t)&=\int_0^x\log(2\cos2\pi v+2\cos t)\,dv-2J(x),\\
 J(x)&=\int_0^x\log(2\sin\pi v)\,dv.
 \end{aligned}
 \label{eq:circle-11}
\end{equation} The stationary branch is the holomorphic solution of
$\partial_x[b(x,t)+b(x,u)]=0$ through $x=1/4$ at $t=u=0$.
The notation $\operatorname{stat}_x$ means substitution of this
solution; logarithm branches are continued from that point. In particular \begin{equation}
 \tan^2(\pi x)=\frac1{1+\tan^2(t/2)+\tan^2(u/2)},\qquad
 \psi_t=-\frac1\pi\operatorname{arsinh}
                       (\tan(t/2)\cos(u/2)).
 \label{eq:circle-12}
\end{equation} Let \(\Psi=\sum_{vw}\psi(\tau_v,\tau_w)\). The first Fourier term of
the exact crossing traces has phase \begin{equation}
 \Psi(\boldsymbol\tau)+\frac{i}{4\pi}\sum_v n_v\tau_v^2
                                      -i\sum_v\tau_v.
 \label{eq:circle-13}
\end{equation} To locate its geometric branch, vary its last coefficient
continuously from zero to one. For \(0\le c\le1\) solve \begin{equation}
 \Psi_{\tau_v}+\frac{i n_v}{2\pi}\tau_v=ic,
                         \qquad\boldsymbol\tau(0)=0.
 \label{eq:circle-14}
\end{equation} The quantitative estimate \eqref{eq:circle-5} gives a unique branch with \begin{equation}
 |\tau_v(c)|\le
 \frac{2\pi c}{(1-1/50)\sqrt{d_v^2+n_v^2}}<0.00111,
                    \qquad\Im\tau_v(c)<0\quad(c>0).
 \label{eq:circle-15}
\end{equation} This is a bound on each local angular parameter. It is not a bound
requiring the total primitive meridian logarithm to stay near zero.

Each edge block admits eight positively oriented ideal tetrahedra
whenever both angular parameters have modulus at most \(0.01\). Their
complete shapes are \(i\), and their deformed shapes stay within
distance \(0.4\) of \(i\). On a disk page put \(Q=e^{i\tau_v}\). The
three complex ratios on its incoming edges are \((-Q^{-1},-Q,-Q^{-1})\);
the outgoing ratios are their inverses. These complex ratios determine the relative positions of the two
triangles, including the displacement along each edge, and depend only on this circle's parameter. Consequently adjacent
blocks fit even if their other endpoints have different parameters. At
the distinguished odd seam the original cross-copy disk maps exchange
the prescribed punctures. Their full matrix edge words are identities.
Positivity and continuation fix every interior edge angle at \(2\pi\).

We must use the original peripheral paths, rather than infer a filling
slope from the exponentiated equations. The meridian crosses
the consecutive disk-page sectors once. A
pure edge contributes \(-2\pi\psi_{\tau_v}\); the single retained half
twist contributes \(-i\epsilon_v\tau_v\). The original primitive
logarithms are therefore \begin{equation}
 U_v=-2\pi\Psi_{\tau_v}-i\epsilon_v\tau_v,\qquad
 V_v=2i\tau_v,
 \qquad U_v+p_vV_v=-2\pi ic.
 \label{eq:circle-16}
\end{equation} The logarithms are continued along the specified peripheral paths.
They are similarity logarithms: after sending the chosen peripheral
fixed point to infinity, they are logarithms of the affine dilation
multiplier. They are twice the corresponding eigenvalue logarithms. Their last identity
cannot be obtained by reducing \(U_v\) to a principal logarithm.

For completeness, the link of an auxiliary ideal vertex is a torus with
a nonsingular similarity structure. Since \(\Re V_v>0\), its holonomy
has a second fixed point \(q_v\) after the marked point has been sent to
infinity. The developing image avoids \(q_v\): a zero of the invariant
vector field \((z-q_v)\partial_z\) has positive index, whereas a torus
has total index zero. Hence \(\log(z-q_v)\) develops a complete
Euclidean torus. Near \(c=0\), \(q_v\) tends to infinity and the periods
on the fixed original paths tend to zero. Their difference from \eqref{eq:circle-16} is
a continuous integer multiple of \(2\pi i\), and is therefore zero along
the whole path.

The slope in \eqref{eq:circle-16} has cone angle \(2\pi c\). Its independent longitude
has positive real period, so it bounds a punctured metric tube with a
positive core length. At \(c=1\), attaching the core produces a smooth
solid torus on precisely the original primitive slope. The main cusp
stays a complete parabolic cusp. This constructs a complete
finite-volume hyperbolic metric on the original knot complement,
including the short coefficients permitted by \eqref{eq:circle-5}. Sections~\ref{sec:fa} and~\ref{sec:fg} give the vertex and face
formulas and the metric completion, including the odd seams and the
continued logarithmic periods.

\section{The action and the invariant have the same
normalization}\label{sec:circle-the-action-and-the-invariant-have-the-same-normalization}

The cells also fix the additive constant of the action. Differential
identities alone would leave that constant undetermined. With
\(V_v=2i\tau_v\), set \begin{equation}
 G=2\pi i\left(\Psi-\sum_v\tau_v\Psi_{\tau_v}\right)
                                  +\frac12\sum_v\epsilon_v\tau_v^2.
 \label{eq:circle-17}
\end{equation} Then \(dG=\frac12\sum_vV_v\,dU_v\). At zero, its imaginary part is
the volume of the original complete augmented link. The marked boundary
formula is \begin{equation}
 \operatorname{Vol}=\Im G-\frac12\sum_v\Re U_v\Im V_v.
 \label{eq:circle-18}
\end{equation} Let \(\mathcal C(c)\) be \eqref{eq:circle-13} with the last term multiplied by \(c\),
evaluated at \eqref{eq:circle-14}. Substitution of the original filling relation in
\eqref{eq:circle-17}--\eqref{eq:circle-18} gives \begin{equation}
                         \operatorname{Vol}_c=2\pi\Re\mathcal C(c).
 \label{eq:circle-19}
\end{equation} At \(c=1\), the right side is thus the volume of the knot complement
just constructed.

It remains to show that this action contributes to the invariant.
Expand the original finite contraction by assigning an integer label
\(0\le k<N\) to every internal edge. Its
normalized label is \(x=k/N\); the high range is
\(k>\lfloor(N-1)/2\rfloor\). The complete contribution of an edge is
maximized at \(x=1/4\). Outside \([1/5,3/10]\), that edge loses a fixed
exponential amount; a high index loses a fixed amount as well. The total
discarded source has an absolute gap \(\pi/200\), independently of graph
size, although its polynomial prefactor can depend on the fixed tree and
powers. The exact finite angular sum has a Gaussian bound weighted by
the vertex degrees. Its real decay constant can be taken as \(1/400\)
and its imaginary growth constant as \(70\).

Preconditioning \eqref{eq:circle-14} by \((d_v-i n_v)/(2\pi)\) gives a contraction with
constant at most \(1/25\). The same estimates compare every retained
nonzero integer Fourier index, also called a charge, with charge one, with a strict loss unless every
charge has magnitude one. The exact crossing-trace identity has already
paired the reflected positive and negative modes, and the resulting
coefficient of the first charge is exactly one. Exact contour identities
bound the omitted charge tails and their joining segments by an absolute
exponential error. The possible drop of the first action is at most
\(\pi(2550/2401)\mathfrak S\), which \eqref{eq:circle-5} places below half the available
error gap. Thus the selected term is nonzero for both parities of \(N\).

After the source and contour errors have been removed, the physical
closing trace contributes exactly \(H_N\), and all
remaining crossing, changes of tensor order and zero-framing factors
combine into the exact \(A_N\) in \eqref{eq:circle-4}. The bounded insertion on each
retained edge is \(1+O(N^{-1})\); the conditioned saddle evaluation
therefore gives \eqref{eq:circle-7}. The full constant and the scalar are retained when
passing to \eqref{eq:circle-6}--\eqref{eq:circle-8}. Neither the volume of a comparison diagram nor the
complete augmented volume replaces the filled knot volume. Appendix~\ref{app:quantitative-sums} gives the finite-sum
proof, including the contraction, Fourier tails and connecting
contours. Sections~\ref{sec:qu-1}--\ref{sec:qu-2} prove the
Newton polynomial bounds and decay constants.
The exact scalar calculation in
\cite[Proposition~8.1]{MeridianTangles} retains the physical strand
orientations and framing in \eqref{eq:circle-4}.

\section{A geometric form of the numerical
condition}\label{sec:circle-a-geometric-form-of-the-numerical-condition}

The numerical condition has a direct interpretation on the original
complete cusps. At an auxiliary cusp the primitive shape is
\(2i c_v/(1-i\epsilon_v c_v)\), where \(c_v=1/d_v\). The normalized
length of a slope means its Euclidean length divided by the square root
of the cusp area. For the primitive filling slope this gives \begin{equation}
 \mathcal L_v^2=\frac{d_v^2+n_v^2}{2d_v},\qquad
                 \mathfrak S=\frac12\sum_v\mathcal L_v^{-2}.
 \label{eq:circle-20}
\end{equation} Thus \eqref{eq:circle-5} is equivalently \(\sum_v\mathcal L_v^{-2}\le2^{-24}\). This
identity explains why high degree can compensate for a small crossing
power. The proof above establishes the completion for these particular
ordered exteriors; it does not import an unrelated normalized-length
criterion or classify all short fillings.

\section{The preferred cusp and the first
correction}\label{sec:circle-the-preferred-cusp-and-the-first-correction}

The first correction records the remaining knot cusp. Normalize its
meridian translation to one. The preferred longitude, characterized
by linking number zero with the knot, then translates by the preferred
cusp shape \(\tau_K\). This geometric marking is fixed before we
compare it with the quantum correction.

Let \(x_e^*\) be the branch in \eqref{eq:circle-12} evaluated at \(\boldsymbol\tau(1)\).
The finite edge-insertion calculation retains one further term: \begin{equation}
 \frac{F_N(q_N^s)}{F_N(1)A_N(s)H_N(s)}
 =1-\frac{\pi s^2}{N}\sum_e\cot(\pi x_e^*)
                 +O_{\mathcal T,\mathbf n,R}(N^{-2}).
 \label{eq:circle-21}
\end{equation} The finite insertion calculation of
\cite[Theorem~1.2 and Section~6.6]{CircleTreeVolume}
has an integrand independent of \(s\), multiplied by
\(1-\pi s^2\sum_e\cot(\pi x_e)/N+O(N^{-2})\), uniformly with derivatives.
At the nondegenerate joint saddle, the normalized expectation of this
coefficient differs from its critical value by \(O(N^{-1})\).
Its explicit factor \(1/N\) makes that an \(O(N^{-2})\) error.
The common saddle corrections cancel in the ratio, and the other
charges and contour terms are exponentially smaller by the quantitative
estimate. This proves the remainder in \eqref{eq:circle-21} under the
present summed criterion.

There is one additional geometric input in interpreting this
coefficient as a preferred cusp shape. The direct peripheral-word
calculation in
\cite[Section~5]{PreferredLongitudes}
constructs the main boundary path from the face maps, follows every
strand once, and computes its meridional homology. Correcting that path
by the resulting integer framing gives, on the same filled branch,
\begin{equation}
                   \tau_K=a(D)+2i\sum_e\cot(\pi x_e^*).
 \label{eq:circle-22}
\end{equation}
This is a classical peripheral-path calculation. It does not use the
quantum correction \eqref{eq:circle-21} to define the cusp shape.
We explain how the marking is carried to the filling point.

Let \(\lambda_{\mathrm{aug}}\) be the preferred main longitude in
the augmented exterior and \(\mu_K\) its meridian. If \(\ell_v\)
is the linking number of the main knot with crossing circle \(v\),
retain the boundary word
\[
 \lambda_{\mathrm{fill}}=\lambda_{\mathrm{aug}}
                \mu_K^{\sum_vp_v\ell_v^2}.
\]
The integers \(p_v\) are fixed throughout the path. Before its
endpoint this is a curve on the unfilled main boundary. Surgery
homology makes it the preferred longitude at the smooth filling.
The peripheral-word calculation supplies both the nonintegral
translation \(2i\sum_e\cot(\pi x_e)\) and the integer correction.
With the crossing and writhe conventions of this chapter, their
combination is exactly \(a(D)=-w(D)-\sum_vn_v\).

The representation used by that calculation agrees with the geometric
cell representation. At augmented completeness, the main meridian
logarithm and the primitive auxiliary longitude logarithms are local
coordinates on its decorated representation space. This is proved in
\cite[Proposition~6.1]{NonparabolicTreeGeometry},
and applied to this particular choice of curves in its Theorem~1.1.
The proof uses the real Lagrangian space of edge and peripheral rows.
In coordinates \(u_j=\log z_j\),
\(v_j=-\log(1-z_j)\), the shape derivative is
\(T_j=z_j/(1-z_j)\), with \(\Im T_j>0\). Pairing a possible
Jacobian-kernel vector with its complex conjugate forces it to vanish.
Thus the specified peripheral logarithms determine the nearby marked
shapes uniquely, including their chosen fixed points.

The exact circle-variable elimination identifies their edge labels by
\eqref{eq:circle-12}. The two translation formulas therefore agree
near augmented completeness. They then continue along the positive
angular path: the main meridian stays nonidentity parabolic, so its
commuting longitude has a unique translation in the frame where the
meridian translates by one. This ratio has no logarithmic integer
ambiguity. The angular face charts remain regular along the path, and
Section~\ref{sec:fg-4} keeps the peripheral logarithms on their
continued branches. Equality consequently holds at the smooth
endpoint, where \(\lambda_{\mathrm{fill}}\) is the preferred
longitude. This proves the identification in \eqref{eq:circle-22}.

Put \(z=2\pi i s\) and \begin{equation}
 \mathcal H_N(z)=\frac{\sinh(z/2)}{N\sinh(z/(2N))},
                         \qquad\mathcal H_N(0)=1.
 \label{eq:circle-23}
\end{equation} Combining the exact factor \eqref{eq:circle-4}, the expansion \eqref{eq:circle-21}, and the geometric
identity \eqref{eq:circle-22} gives \begin{equation}
 \frac{F_N(e^{z/(2N)})}{F_N(1)\mathcal H_N(z)}
   =1+\frac{\tau_K z^2}{8\pi iN}
                         +O_{\mathcal T,\mathbf n,R}(N^{-2}),
                         \qquad |z|\le2\pi R<2\pi.
 \label{eq:circle-24}
\end{equation} All fixed derivatives have the corresponding limits on smaller disks.
The shape in this formula belongs to the filled knot complement, not the
complete augmented link. The exact scalar in \eqref{eq:circle-4} remains part of the
statement; its quadratic expansion is the integer contribution to the
preferred cusp in \eqref{eq:circle-24}.
